\section{Domain adapters}
\label{sec:domains}

The generic kernel becomes useful only when applications are translated into
its interfaces.  This section records robust adapters for the main families of
proofs represented in the repository.

\subsection{Exact algebra and polynomial identities}

Polynomials over integers or rationals are bounded coefficient lists.
Evaluation, addition, multiplication, derivative, and coefficient comparison
are primitive recursive.  A polynomial identity is checked by normalized
coefficient equality, or by a supplied straight-line derivation whose local
ring steps are verified.

For degree-bounded algorithms, traces can carry coefficient-growth bounds and
positive-denominator invariants.  Exact algebra is the easiest part of the
framework because equality is decidable and all objects are finite.  The main
formalization task is a canonical shared code library; parallel encodings
should not be introduced packet by packet.

\subsection{Roots, spectra, and branch data}

For an exact polynomial, stable outputs are:
\begin{itemize}
\item the root multiset approximated by a finite family of disks;
\item clusters with multiplicities;
\item symmetric functions of roots; or
\item individually labeled roots under a separation certificate.
\end{itemize}
A verifier checks disk coverage, disjointness where claimed, residual or
Rouche-type inequalities, and multiplicity counts.  Independent valid
certificates can be compared through a finite Hausdorff or matching estimate,
so \cref{thm:stable-erasure} often avoids cross-precision branch tracking.

For matrices, a rational iteration may return a diagonal or block-diagonal
certificate.  Exact invariants such as trace, determinant, and characteristic
polynomial are stable observations.  Ordered eigenvalues require a proved
cross-certificate comparison; eigenvectors and projectors require spectral
gaps.  At multiplicity, the stable object is the invariant subspace projector
or cluster, not an individual eigenvector.

Exact Puiseux coefficients from arbitrary Cauchy coefficients are not a valid
uniform target.  Exact algebraic coefficient codes plus Newton-polygon and
splitting data are a valid target; otherwise one returns approximate branch
clusters or guarded simple branches.

\subsection{Constructive analysis}

A continuous function on a compact rational box is presented by:
\begin{enumerate}[label=(\roman*)]
\item a fixed dyadic evaluator;
\item a modulus of continuity;
\item a magnitude bound; and
\item any domain or apartness guards needed by partial operations.
\end{enumerate}
Finite nets replace abstract compactness.  The minimum value is approximated by
bounded minimization on a net plus the modulus error.  A minimizer is returned
relationally; a coherent point requires strong convexity, isolation, or another
shrinking estimate.

Integration uses finite rational sums with a partition modulus.  Differentiation
or inverse-function claims require quantitative derivative remainder and
Jacobian guards.  ODE claims are formulated through a finite step method,
local truncation estimates, and a discrete Gronwall accumulator.  The demand
compiler computes how much input precision each step requires.

\subsection{Geometry and topology}

Coordinates are rational or regular names with explicit error budgets.
Incidence, orientation, and side tests are exact only under separation margins;
near a boundary, overlapping branches or relational orientation data are used.
A rotation is stably represented by quaternion or matrix data.  A unit axis is
exported only with a positive gap from the identity or an explicitly supplied
nonzero vector part.

Topological existence arguments are compiled to finite grids, simplicial or
cubical complexes, and parity, Sperner, degree, or chain certificates.  The
verifier checks a finite combinatorial identity.  Refinement relations and
shrinking cells provide approximate points when a coherent location is needed.
The abstract compact space is never manipulated as an infinite object inside
iPRA.

\subsection{Finite probability and statistics}

Start with finite rational laws.  General distributions are accessed through
finite approximations together with tail, truncation, quantization, moment, and
anti-concentration moduli.  Expectations are finite sums.  Concentration bounds
are rational inequalities.  Discontinuous CDF tests at an atom require an
anti-concentration margin or an interval-valued conclusion.

Limit theorems should be stated as finite quantitative approximation theorems:
for requested tolerance \(2^{-n}\), explicit bounds and certificates imply a
test-metric estimate.  A qualitative weak-convergence statement without a rate
is not automatically a primitive-recursive name.

\subsection{Arithmetic geometry and Diophantine finiteness}

Let \(P(p,x)\) be a decidable predicate for coded points and let
\(\height(p,x)\) be a primitive-recursive height.  Assume that all codes of
height at most \(H\) can be enumerated below a primitive-recursive code bound.

\begin{theorem}[Height bound--complete list equivalence]
\label{thm:height-list}
For such a presentation, the following finite data are interconvertible in
iPRA:
\begin{enumerate}[label=(\alph*)]
\item a bound \(H\) with \(P(p,x)\to\height(p,x)\le H\); and
\item a finite list containing every \(x\) satisfying \(P(p,x)\), together
      with a proof of completeness.
\end{enumerate}
The conversions are primitive recursive in the supplied certificate.
\end{theorem}

\begin{proof}
Given \(H\), enumerate every code in the finite height-bounded search region,
filter by \(P\), and remove duplicates by bounded search.  The height bound
proves completeness.  Conversely, from a complete list take the finite maximum
of the heights of its entries.  For any \(x\) satisfying \(P\), completeness
provides a matching list entry, so its height is bounded by the maximum.
\end{proof}

A cardinality bound alone supplies neither datum: a singleton solution set can
have an uncomputably late code.  Thus Faltings- or Siegel-type packets should
separate the fully formalizable bounded enumerator from the deep height-bound
producer.  Mordell--Weil descent similarly separates finite Selmer data from
certified rational representatives on the corresponding coverings.

\subsection{Analytic number theory}

Ineffective constants are handled by certificate interfaces.  A finite
analytic certificate records explicit cutoffs, exceptional alternatives,
zero-free margins, truncation parameters, and rational upper and lower bounds.
The verifier proves that accepted constants imply the requested finite
inequality.  A separate producer theorem establishes uniform availability of
accepted certificates.  This preserves the exact quantifier order and makes
ineffectivity visible rather than embedding it into an unspecified constant.
